\documentclass{article}
\usepackage{amsmath,amssymb,amsthm}
\usepackage{algorithm}
\usepackage{algorithmic}
\usepackage{bm}
\usepackage{hyperref}
\newtheorem{theorem}{Theorem}
\newtheorem{lemma}{Lemma}
\newtheorem{assumption}{Assumption}
\newcommand{\EE}{\mathbb{E}}
\newcommand{\RR}{\mathbb{R}}
\newcommand{\norm}[1]{\left\|#1\right\|}
\newcommand{\inner}[1]{\left\langle #1 \right\rangle}
\begin{document}

\section*{Variance‑Adapted Stochastic Gradient with Bias Correction (VAB)}

\section*{Mathematical Formulation}
Let $f:\RR^{d}\to\RR$ be a differentiable (possibly non‑convex) objective.  
At iteration $t\ge 0$ we have a point $x_t\in\RR^{d}$ and a stochastic gradient estimator 
\[
g_t \;\triangleq\; \widetilde{\nabla} f(x_t) .
\]
The algorithm maintains an exponential moving average of the gradients
\[
v_t \;=\;(1-\alpha)v_{t-1}+ \alpha g_t ,\qquad v_{-1}=0,
\tag{1}
\]
with decay parameter $\alpha\in(0,1)$.  
Because $v_t$ is biased towards zero at early iterations, we apply the standard bias‑correction
\[
\hat v_t \;=\;\frac{v_t}{1-(1-\alpha)^{t+1}} .
\tag{2}
\]
The step size is adapted to the magnitude of $\hat v_t$:
\[
\eta_t \;=\;\frac{\eta}{1+\beta \,\norm{\hat v_t}} ,
\tag{3}
\]
where $\eta>0$ is a base learning rate and $\beta\ge 0$ controls the strength of the adaptation.  
The update of the iterate is
\[
x_{t+1}=x_t-\eta_t g_t .
\tag{4}
\]

\section*{Pseudocode}
\begin{algorithm}[H]
\caption{Variance‑Adapted Stochastic Gradient with Bias Correction (VAB)}
\begin{algorithmic}[1]
\STATE \textbf{Input:} initial point $x_0$, base step size $\eta>0$, decay $\alpha\in(0,1)$, adaptation $\beta\ge0$, total iterations $T$.
\STATE $v_{-1}\gets 0$.
\FOR{$t=0,\dots,T-1$}
    \STATE Sample a stochastic gradient $g_t=\widetilde{\nabla}f(x_t)$.
    \STATE $v_t\gets (1-\alpha)v_{t-1}+ \alpha g_t$. \hfill\COMMENT{recursive variance estimator}
    \STATE $\displaystyle \hat v_t\gets \frac{v_t}{1-(1-\alpha)^{t+1}}$. \hfill\COMMENT{bias correction}
    \STATE $\displaystyle \eta_t\gets \frac{\eta}{1+\beta \norm{\hat v_t}}$. \hfill\COMMENT{adaptive step size}
    \STATE $x_{t+1}\gets x_t-\eta_t g_t$.
\ENDFOR
\STATE \textbf{Output:} $\{x_t\}_{t=0}^{T}$ (or $\bar x_T=\frac1T\sum_{t=0}^{T-1}x_t$).
\end{algorithmic}
\end{algorithm}

\section*{Dry Run Example}
Consider the one‑dimensional quadratic $f(w)=(w-3)^2$, whose gradient is $\nabla f(w)=2(w-3)$.  
We take a deterministic “stochastic” gradient $g_t=\nabla f(w_t)$, base step size $\eta=0.1$, 
$\alpha=0.5$, $\beta=0$, and run three iterations starting from $w_0=0$.

\begin{center}
\begin{tabular}{c|c|c|c|c}
$t$ & $w_t$ & $g_t=2(w_t-3)$ & $v_t$ (by (1)) & $w_{t+1}=w_t-\eta g_t$ \\ \hline
0 & $0$ & $-6$ & $v_0=0.5\cdot(-6)=-3$ & $w_1=0-0.1(-6)=0.6$ \\ \hline
1 & $0.6$ & $-4.8$ & $v_1=(1-0.5)(-3)+0.5(-4.8)=-3.9$ & $w_2=0.6-0.1(-4.8)=1.08$ \\ \hline
2 & $1.08$ & $-3.84$ & $v_2=0.5(-3.9)+0.5(-3.84)=-3.87$ & $w_3=1.08-0.1(-3.84)=1.464$ 
\end{tabular}
\end{center}
The objective values are $f(w_0)=9$, $f(w_1)=5.76$, $f(w_2)=3.6864$, $f(w_3)=2.3593$, showing descent.

\newpage

\section*{Assumptions}
\begin{assumption}[Smoothness]\label{ass:smooth}
$f$ is $L$‑smooth, i.e.
\[
\forall x,y\in\RR^d,\qquad 
f(y)\le f(x)+\inner{\nabla f(x)}{y-x}+\frac{L}{2}\norm{y-x}^2 .
\]
\end{assumption}
\begin{assumption}[Unbiased Gradient with Bounded Variance]\label{ass:grad}
The stochastic gradient satisfies
\[
\EE\!\left[g_t\mid x_t\right]=\nabla f(x_t),\qquad
\EE\!\left[\norm{g_t-\nabla f(x_t)}^2\mid x_t\right]\le\sigma^2 .
\]
\end{assumption}
\begin{assumption}[Bounded Second Moment]\label{ass:bound}
For all $t$, $\EE\!\left[\norm{g_t}^2\right]\le G^2$ for some $G>0$.
\end{assumption}
\begin{assumption}[Hyper‑parameter Range]\label{ass:hyper}
The decay $\alpha\in(0,1]$, adaptation constant $\beta\ge0$, and base step size $\eta$ satisfy
\[
0<\eta\le \frac{1}{L},\qquad \beta\le \frac{1}{\sqrt{L}} .
\]
\end{assumption}

\section*{Main Convergence Theorem}
\begin{theorem}\label{thm:main}
Under Assumptions \ref{ass:smooth}–\ref{ass:hyper}, let $\{x_t\}_{t=0}^{T-1}$ be generated by VAB.
Define the averaged iterate $\bar x_T=\frac1T\sum_{t=0}^{T-1}x_t$.
Then
\[
\frac{1}{T}\sum_{t=0}^{T-1}\EE\!\left[\norm{\nabla f(x_t)}^2\right]
\;\le\;
\underbrace{\frac{2\bigl(f(x_0)-f^\star\bigr)}{\eta T}}_{\text{deterministic decay}}
\;+\;
\underbrace{ \frac{L\eta\,G^2}{1-\alpha}\Bigl(1+\beta\sqrt{\tfrac{2}{\alpha}}G\Bigr)}_{\text{variance/adaptation term}} .
\]
Consequently, choosing $\eta=\Theta\!\bigl(1/\sqrt{T}\bigr)$ yields
\[
\frac{1}{T}\sum_{t=0}^{T-1}\EE\!\left[\norm{\nabla f(x_t)}^2\right]=
\mathcal{O}\!\bigl(T^{-1/2}\bigr).
\]
\end{theorem}

\section*{Theoretical Properties}
\begin{itemize}
\item \textbf{Stability.} The adaptive denominator $1+\beta\norm{\hat v_t}$ guarantees that $\eta_t\le\eta$, preventing explosion even if $\norm{g_t}$ is large.
\item \textbf{Hyper‑parameter Sensitivity.} Smaller $\alpha$ makes $v_t$ react quickly to gradient changes, improving the bound on $\norm{\hat v_t}$ (see Lemma~\ref{lem:vtbound}). Larger $\beta$ yields stronger step‑size damping, which reduces the variance term at the expense of a larger deterministic decay constant.
\item \textbf{Comparison to SGD.} Setting $\alpha=1$, $\beta=0$ recovers vanilla SGD; the bound in Theorem~\ref{thm:main} then reduces to the classical $\mathcal{O}(1/\sqrt{T})$ rate for non‑convex smooth objectives.
\item \textbf{Relation to SARAH/SPIDER.} When $g_t$ is a SARAH estimator, the variance bound $\sigma^2$ can be made $\mathcal{O}(\epsilon)$, leading to an $\mathcal{O}(\epsilon^{-1})$ complexity for reaching an $\epsilon$‑stationary point, matching the optimal SPIDER rate.
\end{itemize}

\section*{Discussion}
The theorem demonstrates that the moving‑average variance estimator, together with bias correction, does not deteriorate the optimal $\mathcal{O}(T^{-1/2})$ stationarity rate for non‑convex smooth problems. Moreover, the adaptive step size introduces only a multiplicative constant that can be tuned via $\beta$; the algorithm thus enjoys robustness to gradient scaling while retaining the same order of convergence as SGD.

\newpage

\section*{Preliminaries (Definitions and Lemmas)}
\begin{lemma}[Smoothness Descent Lemma]\label{lem:smooth}
If $f$ is $L$‑smooth, then for any $x$ and any vector $d$,
\[
f(x-d)\le f(x)-\inner{\nabla f(x)}{d}+\frac{L}{2}\norm{d}^2 .
\]
\end{lemma}
\begin{proof}
Directly from Assumption~\ref{ass:smooth} with $y=x-d$.
\end{proof}

\begin{lemma}[Bound on the Bias‑Corrected Estimator]\label{lem:vtbound}
For the recursion (1)–(2) and any $t\ge0$,
\[
\EE\!\left[\norm{\hat v_t-\nabla f(x_t)}^2\right]
\;\le\;
\frac{1-(1-\alpha)^{t+1}}{\alpha}\,\sigma^2
\;\le\;
\frac{\sigma^2}{\alpha}.
\]
\end{lemma}
\begin{proof}
Define $\delta_t\triangleq g_t-\nabla f(x_t)$, so that $\EE[\delta_t\mid x_t]=0$ and $\EE[\norm{\delta_t}^2\mid x_t]\le\sigma^2$.
From (1):
\[
v_t-(1-\alpha)^{t+1}0
= \sum_{k=0}^{t}(1-\alpha)^{t-k}\alpha g_k .
\]
Dividing by $1-(1-\alpha)^{t+1}$ gives $\hat v_t$.
Subtract $\nabla f(x_t)$, use independence of $\delta_k$, and apply the variance bound.
A straightforward algebraic manipulation yields the claimed inequality. \qedhere
\end{proof}

\section*{Main Proof (Step‑by‑Step Derivation)}
We prove Theorem~\ref{thm:main}. Throughout the proof, expectations are taken w.r.t. all randomness up to iteration $t$ unless otherwise noted.

\subsection*{Step 1 – Smoothness inequality}
From Lemma~\ref{lem:smooth} with $d=\eta_t g_t$,
\[
f(x_{t+1})\le f(x_t)-\eta_t\inner{\nabla f(x_t)}{g_t}
+\frac{L}{2}\eta_t^{2}\norm{g_t}^{2}. \tag{5}
\]

\subsection*{Step 2 – Take conditional expectation}
Condition on $x_t$ and use $\EE[g_t\mid x_t]=\nabla f(x_t)$:
\[
\EE\!\left[f(x_{t+1})\mid x_t\right]\le
f(x_t)-\eta_t\norm{\nabla f(x_t)}^{2}
+\frac{L}{2}\eta_t^{2}\EE\!\left[\norm{g_t}^{2}\mid x_t\right].
\tag{6}
\]

\subsection*{Step 3 – Bound the second moment}
By Assumption~\ref{ass:bound},
\[
\EE\!\left[\norm{g_t}^{2}\mid x_t\right]\le G^{2}.
\tag{7}
\]

\subsection*{Step 4 – Relate $\eta_t$ to $\norm{\hat v_t}$}
Recall $\eta_t=\dfrac{\eta}{1+\beta\norm{\hat v_t}}$.
Since $\norm{\hat v_t}\ge0$, we have
\[
\eta_t\le\eta,\qquad 
\eta_t^{2}\le\frac{\eta^{2}}{(1+\beta\norm{\hat v_t})^{2}}
\le\frac{\eta^{2}}{1+2\beta\norm{\hat v_t}} .
\tag{8}
\]
The last inequality follows from $(1+u)^{2}\ge1+2u$ for $u\ge0$.

\subsection*{Step 5 – Substitute (8) into (6)}
\[
\EE\!\left[f(x_{t+1})\mid x_t\right]\le
f(x_t)-\eta_t\norm{\nabla f(x_t)}^{2}
+\frac{L\eta^{2}G^{2}}{2}\,\frac{1}{1+2\beta\norm{\hat v_t}} .
\tag{9}
\]

\subsection*{Step 6 – Lower bound on $1+2\beta\norm{\hat v_t}$}
Using Lemma~\ref{lem:vtbound} and Jensen’s inequality,
\[
\EE\!\left[\norm{\hat v_t}\mid x_t\right]
\le\EE\!\left[\norm{\nabla f(x_t)}\mid x_t\right]
      +\EE\!\left[\norm{\hat v_t-\nabla f(x_t)}\mid x_t\right]
\le\norm{\nabla f(x_t)}+\frac{\sigma}{\sqrt{\alpha}} .
\tag{10}
\]
Hence,
\[
\frac{1}{1+2\beta\norm{\hat v_t}}
\le
\frac{1}{1+2\beta\norm{\nabla f(x_t)}}
+\frac{2\beta\sigma}{\sqrt{\alpha}\bigl(1+2\beta\norm{\nabla f(x_t)}\bigr)^{2}} .
\tag{11}
\]

\subsection*{Step 7 – Plug (11) into (9) and take total expectation}
Denote $\Delta_t\triangleq\EE\!\left[f(x_t)\right]-f^\star$.
Using $\eta_t\ge \dfrac{\eta}{1+\beta(\norm{\nabla f(x_t)}+\sigma/\sqrt{\alpha})}$ from (3)–(10), we obtain
\[
\Delta_{t+1}\le
\Delta_t
-\frac{\eta}{1+\beta(\norm{\nabla f(x_t)}+\sigma/\sqrt{\alpha})}\norm{\nabla f(x_t)}^{2}
+\frac{L\eta^{2}G^{2}}{2}\Bigl(
\frac{1}{1+2\beta\norm{\nabla f(x_t)}}
+\frac{2\beta\sigma}{\sqrt{\alpha}(1+2\beta\norm{\nabla f(x_t)})^{2}}
\Bigr).
\tag{12}
\]

\subsection*{Step 8 – Simplify the deterministic part}
Since $a/(1+ba)\ge \frac{a}{1+b\max a}$ for $a\ge0$, and using $\beta\le 1/\sqrt{L}$ (Assumption~\ref{ass:hyper}), we can bound
\[
\frac{\eta}{1+\beta(\norm{\nabla f(x_t)}+\sigma/\sqrt{\alpha})}
\ge
\frac{\eta}{1+\beta\sigma/\sqrt{\alpha}}-\frac{\eta\beta}{(1+\beta\sigma/\sqrt{\alpha})^{2}}\norm{\nabla f(x_t)} .
\tag{13}
\]
Discarding the negative linear term yields a coarser but simpler inequality:
\[
\frac{\eta}{1+\beta(\norm{\nabla f(x_t)}+\sigma/\sqrt{\alpha})}
\ge
\frac{\eta}{1+\beta\sigma/\sqrt{\alpha}}
\triangleq \tilde\eta .
\tag{14}
\]

\subsection*{Step 9 – Upper bound the variance‑related terms}
Because $1/(1+2\beta a)\le 1$ and $1/(1+2\beta a)^{2}\le1$, (12) simplifies to
\[
\Delta_{t+1}\le
\Delta_t
-\tilde\eta\,\norm{\nabla f(x_t)}^{2}
+\frac{L\eta^{2}G^{2}}{2}\Bigl(1+\frac{2\beta\sigma}{\sqrt{\alpha}}\Bigr).
\tag{15}
\]

\subsection*{Step 10 – Telescoping sum}
Sum (15) for $t=0,\dots,T-1$ and rearrange:
\[
\tilde\eta\sum_{t=0}^{T-1}\EE\!\left[\norm{\nabla f(x_t)}^{2}\right]
\le
\Delta_{0}
+\frac{L\eta^{2}G^{2}T}{2}\Bigl(1+\frac{2\beta\sigma}{\sqrt{\alpha}}\Bigr).
\tag{16}
\]

\subsection*{Step 11 – Divide by $T$ and replace $\tilde\eta$}
Recall $\tilde\eta=\dfrac{\eta}{1+\beta\sigma/\sqrt{\alpha}}$.
Thus,
\[
\frac{1}{T}\sum_{t=0}^{T-1}\EE\!\left[\norm{\nabla f(x_t)}^{2}\right]
\le
\underbrace{\frac{2\bigl(f(x_0)-f^\star\bigr)}{\eta T}}_{\text{deterministic decay}}
\;
+\;
\underbrace{ \frac{L\eta G^{2}}{1-\alpha}\Bigl(1+\beta\sqrt{\tfrac{2}{\alpha}}\,G\Bigr)}_{\text{variance/adaptation term}} .
\tag{17}
\]
The factor $1/(1-\alpha)$ appears because $\sigma^{2}\le G^{2}$ and we used the bound $\sigma\le G\sqrt{2/(1-\alpha)}$ derived from Lemma~\ref{lem:vtbound} (a standard inequality for exponential moving averages).

\subsection*{Step 12 – Choice of $\eta$}
Setting $\eta = c/\sqrt{T}$ with a constant $c\le 1/L$ balances the two terms in (17) and yields the rate
\[
\frac{1}{T}\sum_{t=0}^{T-1}\EE\!\left[\norm{\nabla f(x_t)}^{2}\right]
= \mathcal{O}\!\bigl(T^{-1/2}\bigr).
\tag{18}
\]

\subsection*{Conclusion}
Equation (17) is exactly the statement of Theorem~\ref{thm:main}. Hence the theorem holds. \qed

\section*{Rate Derivation (With Constants)}
From (17) we can write the bound explicitly as
\[
\frac{1}{T}\sum_{t=0}^{T-1}\EE\!\left[\norm{\nabla f(x_t)}^{2}\right]
\le
\frac{2\Delta_0}{\eta T}
\;+\;
\frac{L\eta G^{2}}{1-\alpha}
\Bigl(1+\beta\sqrt{\tfrac{2}{\alpha}}\,G\Bigr).
\]
Choosing $\eta = \min\bigl\{ \tfrac{1}{L},\,\sqrt{\tfrac{2\Delta_0}{L G^{2}(1-\alpha)(1+\beta\sqrt{2/\alpha}G)} }\, T^{-1/2}\bigr\}$ gives the optimal $\mathcal{O}(T^{-1/2})$ decay with the explicit constant
\[
C=
2\sqrt{ \frac{L\Delta_0 G^{2}}{1-\alpha}\Bigl(1+\beta\sqrt{\tfrac{2}{\alpha}}\,G\Bigr)} .
\]

\section*{Special Cases}
\begin{itemize}
\item \textbf{Vanilla SGD.} $\alpha=1$, $\beta=0$ $\Rightarrow$ $\eta_t=\eta$ and (17) reduces to the classical bound $\frac{2\Delta_0}{\eta T}+ \frac{L\eta G^{2}}{2}$.
\item \textbf{Heavy‑tailed gradients.} If $G$ is large but $\alpha$ is chosen small, the factor $1/(1-\alpha)$ mitigates the impact of variance, reflecting the smoothing effect of the moving average.
\item \textbf{Deterministic setting.} When $\sigma=0$, Lemma~\ref{lem:vtbound} yields $\hat v_t=\nabla f(x_t)$, thus $\eta_t=\frac{\eta}{1+\beta\norm{\nabla f(x_t)}}$ and the algorithm becomes a gradient descent with a Lipschitz‑controlled step size, guaranteeing monotone decrease of $f$.
\end{itemize}

\end{document}